% Options for packages loaded elsewhere
% Options for packages loaded elsewhere
\PassOptionsToPackage{unicode}{hyperref}
\PassOptionsToPackage{hyphens}{url}
\PassOptionsToPackage{dvipsnames,svgnames,x11names}{xcolor}
%
\documentclass[
  11pt,
  a4paper]{article}
\usepackage{xcolor}
\usepackage[left=1in,top=1in,right=1in,bottom=1in]{geometry}
\usepackage{amsmath,amssymb}
\setcounter{secnumdepth}{-\maxdimen} % remove section numbering
\usepackage{iftex}
\ifPDFTeX
  \usepackage[T1]{fontenc}
  \usepackage[utf8]{inputenc}
  \usepackage{textcomp} % provide euro and other symbols
\else % if luatex or xetex
  \usepackage{unicode-math} % this also loads fontspec
  \defaultfontfeatures{Scale=MatchLowercase}
  \defaultfontfeatures[\rmfamily]{Ligatures=TeX,Scale=1}
\fi
\usepackage{lmodern}
\ifPDFTeX\else
  % xetex/luatex font selection
\fi
% Use upquote if available, for straight quotes in verbatim environments
\IfFileExists{upquote.sty}{\usepackage{upquote}}{}
\IfFileExists{microtype.sty}{% use microtype if available
  \usepackage[]{microtype}
  \UseMicrotypeSet[protrusion]{basicmath} % disable protrusion for tt fonts
}{}
\makeatletter
\@ifundefined{KOMAClassName}{% if non-KOMA class
  \IfFileExists{parskip.sty}{%
    \usepackage{parskip}
  }{% else
    \setlength{\parindent}{0pt}
    \setlength{\parskip}{6pt plus 2pt minus 1pt}}
}{% if KOMA class
  \KOMAoptions{parskip=half}}
\makeatother
% Make \paragraph and \subparagraph free-standing
\makeatletter
\ifx\paragraph\undefined\else
  \let\oldparagraph\paragraph
  \renewcommand{\paragraph}{
    \@ifstar
      \xxxParagraphStar
      \xxxParagraphNoStar
  }
  \newcommand{\xxxParagraphStar}[1]{\oldparagraph*{#1}\mbox{}}
  \newcommand{\xxxParagraphNoStar}[1]{\oldparagraph{#1}\mbox{}}
\fi
\ifx\subparagraph\undefined\else
  \let\oldsubparagraph\subparagraph
  \renewcommand{\subparagraph}{
    \@ifstar
      \xxxSubParagraphStar
      \xxxSubParagraphNoStar
  }
  \newcommand{\xxxSubParagraphStar}[1]{\oldsubparagraph*{#1}\mbox{}}
  \newcommand{\xxxSubParagraphNoStar}[1]{\oldsubparagraph{#1}\mbox{}}
\fi
\makeatother


\usepackage{longtable,booktabs,array}
\newcounter{none} % for unnumbered tables
\usepackage{calc} % for calculating minipage widths
% Correct order of tables after \paragraph or \subparagraph
\usepackage{etoolbox}
\makeatletter
\patchcmd\longtable{\par}{\if@noskipsec\mbox{}\fi\par}{}{}
\makeatother
% Allow footnotes in longtable head/foot
\IfFileExists{footnotehyper.sty}{\usepackage{footnotehyper}}{\usepackage{footnote}}
\makesavenoteenv{longtable}
\usepackage{graphicx}
\makeatletter
\newsavebox\pandoc@box
\newcommand*\pandocbounded[1]{% scales image to fit in text height/width
  \sbox\pandoc@box{#1}%
  \Gscale@div\@tempa{\textheight}{\dimexpr\ht\pandoc@box+\dp\pandoc@box\relax}%
  \Gscale@div\@tempb{\linewidth}{\wd\pandoc@box}%
  \ifdim\@tempb\p@<\@tempa\p@\let\@tempa\@tempb\fi% select the smaller of both
  \ifdim\@tempa\p@<\p@\scalebox{\@tempa}{\usebox\pandoc@box}%
  \else\usebox{\pandoc@box}%
  \fi%
}
% Set default figure placement to htbp
\def\fps@figure{htbp}
\makeatother

\ifLuaTeX
  \usepackage{luacolor}
  \usepackage[soul]{lua-ul}
\else
  \usepackage{soul}
\fi

% definitions for citeproc citations
\NewDocumentCommand\citeproctext{}{}
\NewDocumentCommand\citeproc{mm}{%
  \begingroup\def\citeproctext{#2}\cite{#1}\endgroup}
\makeatletter
 % allow citations to break across lines
 \let\@cite@ofmt\@firstofone
 % avoid brackets around text for \cite:
 \def\@biblabel#1{}
 \def\@cite#1#2{{#1\if@tempswa , #2\fi}}
\makeatother
\newlength{\cslhangindent}
\setlength{\cslhangindent}{1.5em}
\newlength{\csllabelwidth}
\setlength{\csllabelwidth}{3em}
\newenvironment{CSLReferences}[2] % #1 hanging-indent, #2 entry-spacing
 {\begin{list}{}{%
  \setlength{\itemindent}{0pt}
  \setlength{\leftmargin}{0pt}
  \setlength{\parsep}{0pt}
  % turn on hanging indent if param 1 is 1
  \ifodd #1
   \setlength{\leftmargin}{\cslhangindent}
   \setlength{\itemindent}{-1\cslhangindent}
  \fi
  % set entry spacing
  \setlength{\itemsep}{#2\baselineskip}}}
 {\end{list}}
\usepackage{calc}
\newcommand{\CSLBlock}[1]{\hfill\break\parbox[t]{\linewidth}{\strut\ignorespaces#1\strut}}
\newcommand{\CSLLeftMargin}[1]{\parbox[t]{\csllabelwidth}{\strut#1\strut}}
\newcommand{\CSLRightInline}[1]{\parbox[t]{\linewidth - \csllabelwidth}{\strut#1\strut}}
\newcommand{\CSLIndent}[1]{\hspace{\cslhangindent}#1}



\setlength{\emergencystretch}{3em} % prevent overfull lines

\providecommand{\tightlist}{%
  \setlength{\itemsep}{0pt}\setlength{\parskip}{0pt}}



 


\usepackage{iftex}

% % Minion Pro
\usepackage[lf]{MinionPro}
\usepackage{MyriadPro}

% IBM Plex
% \usepackage{fontspec}
% \usepackage{unicode-math}
% \setmainfont{IBM Plex Serif}
% \setsansfont{IBM Plex Sans}
% \setmonofont{IBM Plex Mono}
% \setmathfont{IBM Plex Math Beta 240212}

% STIX Two
% \usepackage{fontspec}
% \setmainfont{Stix Two Text}
% \setmathfont{Stix Two Math}
% \setsansfont[Scale=0.9608140457]{Fira Sans} % 7.2489/7.54454

% Libertinus Serif + Libertinus Math + Source Sans Pro + Source Code Pro
% \ifPDFTeX
% \usepackage[mono=false,osf]{libertine}
% \usepackage[libertine]{newtxmath}
% % \usepackage[scale=1,osf]{sourcesanspro} % 7.1832/7.1832
% % \usepackage[scale=0.8772671637]{sourcecodepro} % 7.1832/7.36934*0.9
% \usepackage[scale=0.9181263229,osf]{roboto} % 7.1832/7.82376
% \usepackage[scale=0.8132241635]{roboto-mono} % 7.1832/7.94969*0.9
% \else
% \usepackage[mono=false,osf]{libertine}
% \setmathfont[Scale=MatchUppercase]{libertinusmath-regular.otf}
% % \usepackage[scale=0.9850312868,osf]{sourcesanspro} % 7.2051/7.31459
% % \usepackage[scale=0.8865281581]{sourcecodepro} % 7.2051/7.31459*0.9
% \usepackage[scale=0.9126225152,osf]{roboto} % 7.2051/7.89494
% \usepackage[scale=0.8213602637]{roboto-mono} % 7.2051/7.89494*0.9
% \fi

% Enable section and paragraph numbering
\setcounter{secnumdepth}{4} % Numbering down to paragraph level
\setcounter{tocdepth}{4} % Include paragraphs in the ToC

%%% CUSTOMIZE SECTION SIZES
\makeatletter
\renewcommand{\section}{%
  \@startsection{section}{1}
    {0pt} 					% indent
    {1.0em plus 0.2em minus 0.2em} 	% beforeskip
    {0.001em plus 0.001em minus 0.001em}    	% afterskip
    {\Large\bfseries\sffamily}%
}
\renewcommand{\subsection}{%
  \@startsection{subsection}{2}
    {0pt} 					% indent
    {1.0em plus 0.2em minus 0.2em} 	% beforeskip
    {0.001em plus 0.001em minus 0.001em}    	% afterskip
    {\large\bfseries\sffamily}%
}
\renewcommand{\subsubsection}{%
  \@startsection{subsubsection}{3}
    {0pt} 					% indent
    {1.0em plus 0.2em minus 0.2em} 	% beforeskip
    {0.001em plus 0.001em minus 0.001em}    	% afterskip
    {\normalsize\bfseries\sffamily}%
}
\renewcommand{\paragraph}{%
  \@startsection{paragraph}{4}
    {0pt} 					% indent
    {1.0em plus 0.2em minus 0.2em} 	% beforeskip
    {0.001em plus 0.001em minus 0.001em}    	% afterskip
    {\normalsize\bfseries\sffamily}%
}
\makeatother

%%% GENERAL PACKAGES
\usepackage{hyphenat} % don't hyphenat titles
\usepackage{graphicx}
\usepackage{tabularx} 
\usepackage{marvosym} % Mail logo
\usepackage{enumitem} % Enumerate items (list)
\usepackage{amsmath} % for enhanced math features
\usepackage[labelfont=bf,justification=justified,singlelinecheck=false]{caption} % Figures caption setup
\usepackage[rightcaption]{sidecap} % figure with sidecaption
\usepackage[]{lineno} % Line numbers
\usepackage{ragged2e} % justify
\usepackage{array}
\usepackage{lscape} % landscape tables
\usepackage{rotating} % landscape tables
\usepackage{bm} % bold math

%%%% ORCID LOGO
\usepackage{academicons}
\definecolor{orcidlogocol}{HTML}{A6CE39}
\usepackage{orcidlink}

%%%% LAYOUT 
\tolerance=200
\emergencystretch=2em
\hyphenpenalty=1000
\hbadness=2000
\widowpenalty=1000 % widow penalty
\clubpenalty=1000 % orphans penalty
\renewcommand{\baselinestretch}{1.3333333333333333333333333333} % line spacing

\renewcommand{\arraystretch}{1.3333}
\makeatletter
\@ifpackageloaded{float}{}{\usepackage{float}}
\floatstyle{plain}
\@ifundefined{c@chapter}{\newfloat{suppfig}{h}{losuppfig}}{\newfloat{suppfig}{h}{losuppfig}[chapter]}
\floatname{suppfig}{Figure S}
\newcommand*\quartosuppfigref[1]{Figure \hyperref[#1]{S\ref{#1}}}
\@ifpackageloaded{caption}{}{\usepackage{caption}}
\DeclareCaptionLabelFormat{quartosuppfigreflabelformat}{#1#2}
\captionsetup[suppfig]{labelformat=quartosuppfigreflabelformat}
\newcommand*\listofsuppfigs{\listof{suppfig}{List of Supplementary Figures}}
\makeatother
\makeatletter
\@ifpackageloaded{float}{}{\usepackage{float}}
\floatstyle{plain}
\@ifundefined{c@chapter}{\newfloat{supptbl}{h}{losupptbl}}{\newfloat{supptbl}{h}{losupptbl}[chapter]}
\floatname{supptbl}{Table S}
\floatstyle{plaintop}
\restylefloat{supptbl}
\newcommand*\quartosupptblref[1]{Table \hyperref[#1]{S\ref{#1}}}
\@ifpackageloaded{caption}{}{\usepackage{caption}}
\DeclareCaptionLabelFormat{quartosupptblreflabelformat}{#1#2}
\captionsetup[supptbl]{labelformat=quartosupptblreflabelformat}
\newcommand*\listofsupptbls{\listof{supptbl}{List of Supplementary Tables}}
\makeatother
\makeatletter
\@ifpackageloaded{caption}{}{\usepackage{caption}}
\AtBeginDocument{%
\ifdefined\contentsname
  \renewcommand*\contentsname{Table of contents}
\else
  \newcommand\contentsname{Table of contents}
\fi
\ifdefined\listfigurename
  \renewcommand*\listfigurename{List of Figures}
\else
  \newcommand\listfigurename{List of Figures}
\fi
\ifdefined\listtablename
  \renewcommand*\listtablename{List of Tables}
\else
  \newcommand\listtablename{List of Tables}
\fi
\ifdefined\figurename
  \renewcommand*\figurename{Figure}
\else
  \newcommand\figurename{Figure}
\fi
\ifdefined\tablename
  \renewcommand*\tablename{Table}
\else
  \newcommand\tablename{Table}
\fi
}
\@ifpackageloaded{float}{}{\usepackage{float}}
\floatstyle{ruled}
\@ifundefined{c@chapter}{\newfloat{codelisting}{h}{lop}}{\newfloat{codelisting}{h}{lop}[chapter]}
\floatname{codelisting}{Listing}
\newcommand*\listoflistings{\listof{codelisting}{List of Listings}}
\makeatother
\makeatletter
\makeatother
\makeatletter
\@ifpackageloaded{caption}{}{\usepackage{caption}}
\@ifpackageloaded{subcaption}{}{\usepackage{subcaption}}
\makeatother
\usepackage{bookmark}
\IfFileExists{xurl.sty}{\usepackage{xurl}}{} % add URL line breaks if available
\urlstyle{same}
\hypersetup{
  pdftitle={Supplementary Material -- Title??},
  colorlinks=true,
  linkcolor={blue},
  filecolor={Maroon},
  citecolor={Blue},
  urlcolor={Blue},
  pdfcreator={LaTeX via pandoc}}


\title{Supplementary Material -- Title??}
\author{}
\date{}
\begin{document}

\begin{titlepage}
% This is a combination of Pandoc templating and LaTeX
% Pandoc templating https://pandoc.org/MANUAL.html#templates
% See the README for help
\begin{titlepage}

\raggedright
% Title and subtitle
{\sffamily\LARGE\bfseries\nohyphens{Supplementary Material --
Title??}}\\[2\baselineskip]


% Authors
% This hairy bit of code is just to get "and" between the last 2 authors.
	
\vspace{0.5\baselineskip}

%%%%%% Affiliations & Corresponding author

\vspace{1\baselineskip}

%%%%%% Project website

\vspace{1\baselineskip}

%%%%%% Abstract

\vfill
\vspace{0.1\textheight}
\end{titlepage}\end{titlepage}

\linenumbers

\renewcommand{\thesection}{S\arabic{section}}
\setcounter{section}{0}

\phantomsection

\section{Selection of experimental SSC
conditions}\label{selection-of-experimental-ssc-conditions}

In the experiment, we aimed to measure an informative and representative
subset of stretch-shortening cycles (SSCs) to investigate the effect of
the SSC parameters --- cycle frequency, FTS and MTC length excursion ---
on the maximally attainable AMPO. Yet, a key question was: Which SSCs
should be selected? To make an informed decision, we performed
preliminary simulations using a Hill-type MTC model with parameter
values from literature for three rat m. gastrocnemius medialis
(\citeproc{ref-van_zandwijk_twitch_1996}{Zandwijk et al., 1996}).

We predicted maximal AMPO across various combinations of cycle frequency
and FTS at five distinct MTC length excursions (2, 4, 6, 8 and 10 mm),
each centred at an average MTC length \textasciitilde3\,mm below the MTC
length yielding maximal isometric GM force. Muscle stimulation onset
time (i.e., the instance at which stimulation switches from 0 to its
maximal value) was set to the start of MTC shortening. Muscle
stimulation duration (i.e., the time period during which stimulation
remains at its maximal value before returning to 0) was optimized using
the Nelder-Mead simplex method (\citeproc{ref-gao_implementing_2012}{Gao
and Han, 2012}). To ensure periodic behaviour, optimisations were run
for at least five cycles and until the muscle force at the beginning of
the cycle deviated less than 20 mN from its value at the end of the
cycle. Maximal AMPO was calculated for last cycle.

We then made contour plots for every MTC length excursion
(\quartosuppfigref{suppfig-s-prelim}), showing the influence of cycle
frequency and FTS on maximal AMPO at each MTC length excursion. For the
experimental SSC conditions, we limited the range of FTS between 0.2 and
0.8 because values outside this range will yield very short shortening
or lengthening durations. Within this FTS range, we observed that AMPO
peaked at about 5 and 3 Hz for an MTC length excursion of 4 and 8\,mm
respectively. Again, to prevent short shortening and lengthening
duration, we set the maximum cycle frequencies accordingly. Based on
these findings, we selected three subsets of SSC conditions: 1) constant
FTS, varying cycle frequency; 2) constant cycle frequency, varying FTS;
3) varying both cycle frequency and FTS. Because there was overlap in
the SSC conditions, this yielded 13 unique SSC conditions for each MTC
length excursion, such that there were 26 different SSC conditions in
total.

\begin{suppfig}

\centering{

\pandocbounded{\includegraphics[keepaspectratio]{../../figures/s_prelim.pdf}}

}

\caption{\label{suppfig-s-prelim}\textbf{Preliminary predictions of
maximal AMPO as a function of cycle frequency and FTS.} The contour
lines depict the maximally attainable AMPO (in mW) averaged over three
parameter sets. These preliminary predictions were derived with a
Hill-type MTC model with parameters obtained from literature
(\citeproc{ref-van_zandwijk_twitch_1996}{Zandwijk et al., 1996}) and
were used to inform the selection of experimental SSC conditions.}

\end{suppfig}%

\newpage{}

\renewcommand{\thesection}{S\arabic{section}}
\setcounter{section}{1}

\phantomsection

\section{Experimental SSC conditions and experimentally measured
AMPO}\label{experimental-ssc-conditions-and-experimentally-measured-ampo}

\begin{supptbl}[H]

\caption{\label{supptbl-sscpa}SSC parameters, stimulation durations, and
measured AMPO of experimental stretch-shortening cycles with a 4 mm MTC
length excursion. Stimulation onset was set at the start of MTC
shortening in all conditions}

\centering{

\centering
\rotatebox{90}{
  \begin{minipage}{\textheight}
    \centering
    \begin{tabular}{|l|c|c|c|c|c|c|c|c|c|c|c|c|c|c|} \hline
  & & \multicolumn{13}{c|}{Condition}  \\ \hline
  \bfseries & \bfseries Unit & \bfseries 1 & \bfseries 2 & \bfseries 3 & \bfseries 4 & \bfseries 5 & \bfseries 6 & \bfseries 7 & \bfseries 8 & \bfseries 9 & \bfseries 10 & \bfseries 11 & \bfseries 12 & \bfseries 13 \\ \hline
  \multicolumn{15}{|l|}{\itshape SSC parameters} \\ \hline
Cycle frequency & Hz & 1.0 & 2.0 & 3.0 & 4.0 & 5.0 & 3.0 & 3.0 & 3.0 & 3.0 & 5.0 & 4.0 & 2.0 & 1.0 \\ \hline
FTS & - & 0.50 & 0.50 & 0.50 & 0.50 & 0.50 & 0.80 & 0.65 & 0.35 & 0.20 & 0.80 & 0.65 & 0.35 & 0.20 \\ \hline
  \multicolumn{15}{|l|}{\itshape MTC shortening and lengthening times} \\ \hline
MTC shortening time & ms & 500 & 250 & 167 & 125 & 100 & 267 & 217 & 117 & 66.7 & 160 & 162 & 175 & 200 \\ \hline
MTC lengthening time & ms & 500 & 250 & 167 & 125 & 100 & 66.7 & 117 & 217 & 267 & 40.0 & 87.5 & 325 & 800 \\ \hline
  \multicolumn{15}{|l|}{\itshape Stimulation durations} \\ \hline
Trial 1 & ms & 405 & 175 & 103 & 65 & 35 & 183 & 133 & 53 & 23 & 85 & 95 & 115 & 135 \\ \hline
Trial 2 & ms & 445\textsuperscript{1} & 205 & 123 & 85 & 55 & 213 & 163 & 73 & 33\textsuperscript{2} & 115 & 115 & 135 & 155\textsuperscript{3} \\ \hline
  \multicolumn{15}{|l|}{\itshape Measured AMPO of Rat 1} \\ \hline
 Trial 1  & mW  & 31  & 51  & 59  & 58  & 45  & 74  & 67  & 38  & 20  & 82  & 74  & 43  & 23 \\ \hline 
 Trial 2  & mW  & 34  & 56  & 69  & 72  & 66  & 84  & 77  & 49  & 25  & 104  & 87  & 47  & 27 \\ \hline 
  \multicolumn{15}{|l|}{\itshape Measured AMPO of Rat 2} \\ \hline
 Trial 1  & mW  & 29  & 48  & 57  & 58  & 45  & 71  & 62  & 37  & 18  & 79  & 70  & 40  & 22 \\ \hline 
 Trial 2  & mW  & 30  & 52  & 65  & 68  & 64  & 78  & 72  & 46  & 24  & 93  & 79  & 45  & 24 \\ \hline 
  \multicolumn{15}{|l|}{\itshape Measured AMPO of Rat 3} \\ \hline
 Trial 1  & mW  & 34  & 56  & 65  & 66  & 48  & 83  & 72  & 44  & 22  & 91  & 82  & 47  & 25 \\ \hline 
 Trial 2  & mW  & 37  & 62  & 74  & 82  & 74  & 91  & 85  & 56  & 20  & 116  & 96  & 54  & 28 \\ \hline 
\end{tabular}\break\hfill\footnotesize{\textsuperscript{1} For conditon 1: stimulation duration of rat 3 was 455 ms. \textsuperscript{2} For conditon 9: stimulation duration of rat 3 was 23 ms. \textsuperscript{3} For conditon 13: stimulation duration of rat 1 was 165 ms.}
  \end{minipage}
}

}

\end{supptbl}%

\begin{supptbl}[H]

\caption{\label{supptbl-sscpb}SSC parameters, stimulation durations, and
measured AMPO of experimental stretch-shortening cycles with a 8 mm MTC
length excursion. Stimulation onset was set at the start of MTC
shortening in all conditions}

\centering{

\centering
\rotatebox{90}{
  \begin{minipage}{\textheight}
    \centering
    \begin{tabular}{|l|c|c|c|c|c|c|c|c|c|c|c|c|c|c|} \hline
  & & \multicolumn{13}{c|}{Condition}  \\ \hline
  \bfseries & \bfseries Unit & \bfseries 1 & \bfseries 2 & \bfseries 3 & \bfseries 4 & \bfseries 5 & \bfseries 6 & \bfseries 7 & \bfseries 8 & \bfseries 9 & \bfseries 10 & \bfseries 11 & \bfseries 12 & \bfseries 13 \\ \hline
  \multicolumn{15}{|l|}{\itshape SSC parameters} \\ \hline
Cycle frequency & Hz & 1.0 & 1.5 & 2.0 & 2.5 & 3.0 & 2.0 & 2.0 & 2.0 & 2.0 & 3.0 & 2.5 & 1.5 & 1.0 \\ \hline
FTS & - & 0.50 & 0.50 & 0.50 & 0.50 & 0.50 & 0.80 & 0.65 & 0.35 & 0.20 & 0.80 & 0.65 & 0.35 & 0.20 \\ \hline
  \multicolumn{15}{|l|}{\itshape MTC shortening and lengthening times} \\ \hline
MTC shortening time & ms & 500 & 333 & 250 & 200 & 167 & 400 & 325 & 175 & 100 & 267 & 260 & 233 & 200 \\ \hline
MTC lengthening time & ms & 500 & 333 & 250 & 200 & 167 & 100.0 & 175 & 325 & 400 & 66.7 & 140 & 433 & 800 \\ \hline
  \multicolumn{15}{|l|}{\itshape Stimulation durations} \\ \hline
Trial 1 & ms & 405 & 253 & 175 & 135 & 103 & 295 & 245 & 115 & 55 & 173 & 175 & 163 & 145 \\ \hline
Trial 2 & ms & 455 & 293 & 205 & 165 & 123 & 345 & 275 & 135 & 75\textsuperscript{1} & 213\textsuperscript{2} & 225 & 193 & 165\textsuperscript{3} \\ \hline
  \multicolumn{15}{|l|}{\itshape Measured AMPO of Rat 1} \\ \hline
 Trial 1  & mW  & 53  & 67  & 75  & 80  & 79  & 94  & 89  & 55  & 26  & 113  & 95  & 54  & 34 \\ \hline 
 Trial 2  & mW  & 54  & 47  & 80  & 86  & 85  & 99  & 97  & 63  & 28  & 121  & 104  & 62  & 35 \\ \hline 
  \multicolumn{15}{|l|}{\itshape Measured AMPO of Rat 2} \\ \hline
 Trial 1  & mW  & 49  & 63  & 73  & 75  & 76  & 88  & 83  & 52  & 24  & 104  & 88  & 50  & 32 \\ \hline 
 Trial 2  & mW  & 50  & 66  & 76  & 80  & 81  & 92  & 86  & 55  & 27  & 113  & 95  & 54  & 33 \\ \hline 
  \multicolumn{15}{|l|}{\itshape Measured AMPO of Rat 3} \\ \hline
 Trial 1  & mW  & 57  & 73  & 81  & 87  & 88  & 103  & 95  & 60  & 28  & 123  & 101  & 60  & 37 \\ \hline 
 Trial 2  & mW  & 47  & -  & 89  & -  & 88  & 110  & -  & -  & 32  & 129  & -  & -  & 40 \\ \hline 
\end{tabular}\break\hfill\footnotesize{\textsuperscript{1} For conditon 1: stimulation duration of rat 3 was 455 ms. \textsuperscript{2} For conditon 9: stimulation duration of rat 3 was 23 ms. \textsuperscript{3} For conditon 13: stimulation duration of rat 1 was 165 ms.}
  \end{minipage}
}

}

\end{supptbl}%

\newpage{}

\renewcommand{\thesection}{S\arabic{section}}
\setcounter{section}{2}

\phantomsection

\section{Optimisation protocol to identify optimal MTC length over
time}\label{optimisation-protocol-to-identify-optimal-mtc-length-over-time}

One of the aims of this study was to identify the optimal MTC length
over time for maximal AMPO, when MTC length over time is entirely free
to take any form. To predict this, we used a Hill-type MTC model. In our
Hill-type MTC model, we assumed that both SEE and PEE were purely
elastic elements (e.g., \citeproc{ref-anderson_dynamic_1999}{Anderson
and Pandy, 1999}; \citeproc{ref-van_soest_contribution_1993}{Soest and
Bobbert, 1993}). Therefore, during fully periodic SSC neither PEE nor
SEE contributed to AMPO of the MTC. This also means that neither PEE nor
SEE influences the optimal CE length over time that yields maximal AMPO
over a full periodic SSC. We used this to simplify the optimisation
problem. We reformulated to model such that we considered a Hill-type
model of CE only. As such, CE length was no longer a state variable, but
an input to the model. The model therefore had two inputs: CE length
(\(L_{CE}\)) and muscle stimulation (\(STIM\)) over time. The only state
variable was \(\gamma\), the normalised intrafilament \(Ca^{2+}\)
concentration. Given this formulation, \(L_{CE}\), \(STIM\) and
\(\gamma\) were known at each time instant during the optimisation
process. Consequently, CE velocity can be calculated by taking the
time-derivative of CE length. Using CE velocity, CE force could then be
computed using the CE force--velocity relationship. MTC length over time
was subsequently obtained by taking the following steps: 1) PEE force
was calculated from PEE length (which equals CE length); 2) SEE force
was calculated by taking the sum of CE and PEE force; 3) SEE length was
calculated from SEE force; 4) MTC length was calculated as the sum of CE
and SEE length. By computing MTC length over time via these steps, we
could also impose MTC length excursion.

The optimisation problem was thus to find the periodic CE length,
\(STIM\) and \(\gamma\) over time for maximal AMPO. Optimisations were
performed for 1) imposed cycle frequencies (1.0--6.0 Hz in 0.5 Hz
steps); 2) imposed FTS values (0.05--0.95 in 0.05 steps) and 3) imposed
MTC length excursions (2--11\,mm in 1\,mm steps). For all optimisation,
the following cost function was maximised (note that, in our
formulation, CE shortening corresponds to negative CE velocity):

\begin{equation}\protect\phantomsection\label{eq-objective}{
J = -\int_{t_0}^{t_0+T_c} F_{CE}^{rel} \cdot v_{CE}^{rel} \ dt
}\end{equation} where \(t_0\) denotes the initial time and \(T_c\)
denotes the cycle time.

The solution had to satisfy the differential equation describing the
\hl{contraction and excitation} \st{activation} dynamics:
\begin{equation}\protect\phantomsection\label{eq-dynamic-con}{
\dot{\gamma} = f(\gamma(t),STIM(t))
}\end{equation}

The following task constraints were added to achieve periodic behaviour:
\begin{equation}\protect\phantomsection\label{eq-task-con}{
  L_{CE}(t_0) = L_{CE}(t_0+T_c) \\
  \gamma(t_0) = \gamma(t_0+T_c)
}\end{equation}

Moreover, inequality constraints were imposed on \(\gamma(t)\),
\(L_{CE}^{rel}(t)\) and \(STIM(t)\):
\begin{equation}\protect\phantomsection\label{eq-ineq-con1}{
  \gamma_0 \leq \gamma(t) \leq 1 
}\end{equation}

\[
  (1-w) \leq L_{CE}^{rel}(t) \leq (1+w) \\
  0 \leq STIM(t) \leq 1
\] where \(L_{CE}^{rel}\) denotes the CE length normalised the optimum
CE length (i.e.~the CE length yielding maximal isometric CE force).

We implemented our problem as a Multiphase Optimal Control Problem to
assure that CE shortening and lengthening occurred only once per full
periodic cycle, which was done by inequality constraints on
\(v_{CE}^{rel}(t)\):
\begin{equation}\protect\phantomsection\label{eq-fts-con}{
\begin{aligned}
  v_{CE}^{rel}(t_0 \ ... \ t_0 + T_c \cdot FTS) &\leq 0 \\
  v_{CE}^{rel}(t_0 + T_c \cdot FTS \ ... \ t_0 + T_c) &\geq 0
\end{aligned}
}\end{equation} FTS was a numerical value (between 0.05 and 0.95) when
FTS were imposed (see above). When FTS was not imposed, FTS was a
control variable optimised for maximal AMPO.

When MTC length excursion was imposed, an extra inequality constraint
was added:

\begin{equation}\protect\phantomsection\label{eq-lcecon}{
  L_{MTC}(t_0) + L_{MTC}(t_0 + T_c \cdot FTS) = \textrm{MTC length excursion}
}\end{equation}

The dynamic, task and (in)equality constraints were implemented in
CasADi (\citeproc{ref-andersson_casadi_2019}{Andersson et al., 2019})
and the cost function was minimised using a Direct Collocation method
using a nonlinear interior point method (IPOPT,
\citeproc{ref-wachter_implementation_2006}{Wächter and Biegler, 2006}).
The optimal control solution was checked by running forward simulations
of the model with a linear-interpolated \(F_{CE}^{rel}(t)\) and
\(STIM(t)\) as inputs to confirm that no relevant constraints were
violated between the collocation points. This was done using SciPy's ODE
integrator with an Implicit Runge-Kutta method of Radau IIA family of
order 5 (\citeproc{ref-hairer_implementation_1996}{Hairer and Wanner,
1996}) and an absolute and relative tolerance of 1e-6 and a maximum
timestep of 1e-3. This yielded only very small differences between the
states at the collocation points and those at the same time instance of
the forward simulation for all optimisations.

\newpage{}

\renewcommand{\thesection}{S\arabic{section}}
\setcounter{section}{3}

\phantomsection

\section{Optimal CE length over time: influence of FTS and MTC length
excursion}\label{optimal-ce-length-over-time-influence-of-fts-and-mtc-length-excursion}

Using the Hill-type MTC model, we identified the optimal CE length over
time for maximal AMPO without imposing any constraints on MTC length
over time (in contrast to section \textbf{?@sec-r-am}, where MTC
shortening and lengthening velocity were constant). When imposing cycle
frequency, the resulting MTC length over time was reasonably similar to
those when with constant MTC shortening and lengthening velocity (see
\textbf{?@sec-r-oc} and \textbf{?@fig-r-oc}). However, when imposing FTS
or MTC length excursion the shape of CE and MTC length over time
dramatically changed.

For imposed FTS values below the optimal value, CE lengthening velocity
was found to be almost zero at the onset of CE lengthening
(\quartosuppfigref{suppfig-s-oc}A). This infinitely low velocity at the
start of CE lengthening was optimal for AMPO because it allowed the
active state to decrease fully before lengthening occurred, thereby
reducing the amount of negative mechanical work performed. On top of
that, it allowed the active state to remain high throughout most part of
CE shortening to enable substantial positive mechanical work. Towards
the end of CE lengthening, CE velocity again approached zero, which made
it possible to initiate muscle stimulation before CE shortening began.
This resulted to a high active state at the onset of CE shortening to
enhance positive mechanical work. For imposed FTS values above the
optimal value, CE shortening velocity decreased over time at the end of
CE shortening (\quartosuppfigref{suppfig-s-oc}A). To avoid substantial
negative mechanical work during the subsequent CE lengthening, muscle
stimulation ceased well before the end of CE shortening. While CE could
theoretically maintain a constant shortening velocity throughout the
shortening phase, any additional positive mechanical work gained near
the end was outweighed by the increased negative mechanical work. Thus,
an almost isometric phase preceding CE lengthening is optimal for
maximising AMPO. In sum, the CE length over time for imposed FTS values
mostly results from the effect of the excitation dynamics.

For imposed MTC length excursions above the optimal value, CE shortening
velocity was found to be higher at the beginning and end of the
shortening phase compared to the middle part
(\quartosuppfigref{suppfig-s-oc}B). At the start and end of the
shortening phase, only limited positive mechanical work could be
produced due to the low active state and the suboptimal CE length (that
is due to the CE force-length relationship). While CE could
theoretically shorten at a constant velocity throughout the entire
shortening phase, this would reduce the positive mechanical work during
the middle part due to the effect of the force-velocity relationship. At
the same time, little additional positive mechanical work could be
achieved at the beginning and end of the shortening phase (due to the
low active state and suboptimal CE length). Consequently, a non-constant
CE shortening velocity --- with CE shortening first decreasing and then
increasing during the shortening phase --- is optimal for maximising
AMPO at imposed MTC length excursions above the optimal value. For
imposed MTC length excursions below the optimal value, the opposite
pattern was found: CE shortening velocity was found to be slower at the
beginning and end of the shortening phase compared to the middle part
(\quartosuppfigref{suppfig-s-oc}B). As explained, only limited positive
mechanical work could be produced at the beginning and end of the
shortening phase due to the low active state. If CE would shorten with a
constant velocity throughout the entire shortening phase, it would
decrease the shortening distance during the middle part (at which the
active state is high) and thereby reducing the positive mechanical work.
As a result, a non-constant CE shortening velocity --- with a slow
shortening velocity at the end and beginning of the shortening phase ---
is optimal for maximising AMPO at imposed MTC length excursions above
the optimal value. In sum, the optimal CE length over time for imposed
MTC length excursions shows the intricate trade-off between the CE
force-length relationship, CE force-velocity relationship and the and
excitation dynamics.

\begin{suppfig}

\centering{

\pandocbounded{\includegraphics[keepaspectratio]{../../figures/s_oc.pdf}}

}

\caption{\label{suppfig-s-oc}\textbf{Predicted CE length over time for
maximally attainable AMPO, shown for two imposed FTS values (A) and two
imposed MTC length excursions (B), at a cycle frequency of 3.5 Hz for
rat 1}. In panel A, only cycle frequency and FTS were imposed, while in
panel B only cycle frequency and MTC length excursions were imposed.
Apart from these two constraints, CE and MTC length over time were
completely unconstrained and thus the shape of CE length over time could
be different between every combination of imposed cycle frequency and
FTS (A) or imposed cycle frequency and MTC length excursion (B). CE
stimulation was maximal during the period indicated by the coloured bars
and off elsewhere.}

\end{suppfig}%

\begin{suppfig}

\centering{

\pandocbounded{\includegraphics[keepaspectratio]{../../figures/s_ampo.pdf}}

}

\caption{\label{suppfig-s-ampo}\textbf{Predicted influence of cycle
frequency (A), FTS (B) and MTC length excursion (C) on the maximum
attainable AMPO}. The maximally attainable AMPO --- averaged across the
three rats --- slightly increased when there was no constraint on MTC
length over time (orange lines), compared to SSCs with a constant MTC
shortening and lengthening velocity (black lines).}

\end{suppfig}%

\protect\phantomsection\label{refs}
\begin{CSLReferences}{1}{1}
\bibitem[\citeproctext]{ref-anderson_dynamic_1999}
\textbf{Anderson, F. C. and Pandy, M. G.} (1999).
\href{https://doi.org/10.1080/10255849908907988}{A {Dynamic}
{Optimization} {Solution} for {Vertical} {Jumping} in {Three}
{Dimensions}}. \emph{Computer Methods in Biomechanics and Biomedical
Engineering} \textbf{2}, 201--231.

\bibitem[\citeproctext]{ref-andersson_casadi_2019}
\textbf{Andersson, J. A. E., Gillis, J., Horn, G., Rawlings, J. B. and
Diehl, M.} (2019).
\href{https://doi.org/10.1007/s12532-018-0139-4}{{CasADi}: A software
framework for nonlinear optimization and optimal control}.
\emph{Mathematical Programming Computation} \textbf{11}, 1--36.

\bibitem[\citeproctext]{ref-gao_implementing_2012}
\textbf{Gao, F. and Han, L.} (2012).
\href{https://doi.org/10.1007/s10589-010-9329-3}{Implementing the
{Nelder}-{Mead} simplex algorithm with~adaptive parameters}.
\emph{Computational Optimization and Applications} \textbf{51},
259--277.

\bibitem[\citeproctext]{ref-hairer_implementation_1996}
\textbf{Hairer, E. and Wanner, G.} (1996).
\href{https://doi.org/10.1007/978-3-642-05221-7_8}{Implementation of
{Implicit} {Runge}-{Kutta} {Methods}}. In \emph{Solving {Ordinary}
{Differential} {Equations} {II}: {Stiff} and {Differential}-{Algebraic}
{Problems}} (ed. Hairer, E.) and Wanner, G.), pp. 118--130. Berlin,
Heidelberg: Springer Berlin Heidelberg.

\bibitem[\citeproctext]{ref-van_soest_contribution_1993}
\textbf{Soest, A. J. van and Bobbert, M. F.} (1993).
\href{https://doi.org/10.1007/BF00198959}{The contribution of muscle
properties in the control of explosive movements}. \emph{Biological
Cybernetics} \textbf{69}, 195--204.

\bibitem[\citeproctext]{ref-wachter_implementation_2006}
\textbf{Wächter, A. and Biegler, L. T.} (2006).
\href{https://doi.org/10.1007/s10107-004-0559-y}{On the implementation
of an interior-point filter line-search algorithm for large-scale
nonlinear programming}. \emph{Mathematical Programming} \textbf{106},
25--57.

\bibitem[\citeproctext]{ref-van_zandwijk_twitch_1996}
\textbf{Zandwijk, J. P. van, Bobbert, M. F., Baan, G. C. and Huijing, P.
A.} (1996). \href{https://doi.org/10.1007/s004220050306}{From twitch to
tetanus: Performance of excitation dynamics optimized for a twitch in
predicting tetanic muscle forces}. \emph{Biological Cybernetics}
\textbf{75}, 409--417.

\end{CSLReferences}




\end{document}
